% Addition to root-owned NN-S016; not a separate registered finding.
% Standard amsmath/amssymb/amsthm environments only.
\begin{proposition}[NN-S016: complex homotopy type and the affine-space case]
Retain the sparse quadratic equations $f_1,\ldots,f_h$ of NN-C027,
with $v\geq1$ and
\[
 Y=V(f_1,\ldots,f_h)\simeq\coprod_N\operatorname{Spec}\mathbb Z,
 \qquad
 X=V\left(t^2-t-\sum_{i=1}^h y_i f_i(x)\right).
\]
The complex realization has homotopy type
\[
                      X(\mathbb C)\simeq\bigvee_{a=1}^{N}S^{2v},
\]
where the empty wedge means a point. In particular its integral
cohomology is $\mathbb Z$ in degree zero, $\mathbb Z^N$ in degree $2v$,
and zero otherwise, with every product of positive-degree classes zero.
Moreover,
\[
 N=0\quad\Longleftrightarrow\quad
 X\simeq\mathbb A^{v+h}_{\mathbb Z}
 \quad\Longleftrightarrow\quad
 X_{\mathbb C}\simeq\mathbb A^{v+h}_{\mathbb C}
 \quad\Longleftrightarrow\quad X(\mathbb C)\text{ is contractible}.
\]
The first isomorphism, when it exists, is proved by an explicit
polynomial formula from any identity $\sum_i a_i f_i=1$.
No Quillen--Suslin theorem or effective bound for such an identity is
required.
\end{proposition}

\begin{proof}
Write $\pi:X(\mathbb C)\to\mathbb C^v$ for the projection, and identify
$Y(\mathbb C)$ with its $N$ distinct points. First suppose $N>0$.
We prove the local models and the attaching maps before determining
the homotopy type.

Put $U=\mathbb C^v\setminus Y(\mathbb C)$ and
$A=\pi^{-1}(U)$. For $x\in U$ define the continuous column vector
\[
 a(x)=\frac{(\overline{f_1(x)},\ldots,\overline{f_h(x)})}
                  {\sum_i|f_i(x)|^2},\qquad f(x)\cdot a(x)=1.
\]
Every point of $A$ can uniquely be written
\[
       (x,t,y),\qquad y=a(x)t(t-1)+w,\quad f(x)\cdot w=0.
\]
Thus $A$ is the total space of the kernel vector bundle over
$U\times\mathbb C_t$. More explicitly, with $0\leq s\leq1$, set
\[
 t_s=(1-s)t,\qquad w_s=(1-s)w,\qquad
 y_s=a(x)t_s(t_s-1)+w_s.
\]
This is a strong deformation retraction onto the section $t=0,y=0$,
which is $U$. It is compatible with restriction to every open subset
of $U$. In particular projection identifies all intersection maps
below with the corresponding base-space inclusions. This is a
continuous construction; conjugation is not asserted to be algebraic.

For each $a\in Y(\mathbb C)$, reducedness says that the germs of the
$f_i$ generate the maximal ideal in the regular analytic local ring
of $\mathbb C^v$ at $a$. Their differentials therefore span the
cotangent space. Choose $v$ of them with independent differentials.
After reordering, the holomorphic inverse function theorem makes
\[
                    u_i=f_i(x),\qquad 1\leq i\leq v,
\]
coordinates on a neighborhood of $a$. In particular $h\geq v$.
Choose that neighborhood to map biholomorphically onto a small
Euclidean ball $B_a=\{\|u\|<\delta_a\}$ centered at zero, with these
neighborhoods pairwise disjoint. Each remaining holomorphic function
vanishes at zero and has an expression
\[
 f_j(x(u))=\sum_{i=1}^v b_{ji}(u)u_i\quad(j>v).
\]
For example the coefficients are obtained by integrating the partial
derivatives of $f_j(x(su))$ from $s=0$ to $s=1$, which is holomorphic
on the ball. The invertible triangular change of fiber coordinates
\[
 \eta_i=y_i+\sum_{j>v}b_{ji}(u)y_j\quad(1\leq i\leq v)
\]
turns the equation into
\[
                 \sum_{i=1}^v u_i\eta_i=t(t-1),
\]
with $h-v$ unused affine coordinates. Consequently
$V_a=\pi^{-1}(B_a)$ is analytically isomorphic to
\[
 \left\{(u,\eta,t):\sum_i u_i\eta_i=t(t-1),\ \|u\|<\delta_a\right\}
                   \times\mathbb C^{h-v}.
\]
Here and below $B_a$ also denotes its inverse image in the original
base coordinates.

The restriction $\|u\|<\delta_a$ does not change the homotopy type of
the full affine quadric $Q_{2v}:\sum u_i\eta_i=t(t-1)$.
Indeed take
\[
 c(u)=\min\left(1,\frac{\delta_a}{2\|u\|}\right),\quad c(0)=1,
 \qquad c_s(u)=1-s+s c(u).
\]
The maps $(u,\eta,t)\mapsto(c_s(u)u,c_s(u)^{-1}\eta,t)$ stay on the
quadric. At $s=1$ their image has $\|u\|<\delta_a$, and, if it starts
in that open subset, it remains there throughout the homotopy.
They supply a homotopy inverse to the inclusion of the restricted
quadric. This is a homotopy equivalence, not a claim that the whole
restricted subset is pointwise fixed during the homotopy.

The full quadric is linearly equivalent over $\mathbb C$ to
$\sum_{j=0}^{2v}z_j^2=1$: use $z_0=2t-1$ and, for each $i$,
$z_{2i-1}=\mathrm i(u_i+\eta_i)$, $z_{2i}=u_i-\eta_i$.
Writing $z=p+\mathrm i q$ with real vectors gives
\[
 |p|^2-|q|^2=1,\qquad p\cdot q=0.
\]
It is therefore diffeomorphic to the tangent bundle of $S^{2v}$,
by $(e,q)\mapsto\sqrt{1+|q|^2}\,e+\mathrm i q$ with
$|e|=1$ and $e\cdot q=0$. Scaling $q$ to zero retracts this bundle
onto its zero section. We have proved $V_a\simeq S^{2v}$, including
simple connectivity. On the intersection
\[
             W_a=A\cap V_a=\pi^{-1}(B_a\setminus\{a\}),
\]
the earlier projection retraction gives $W_a\simeq S^{2v-1}$.

Put $d=2v\geq2$. The small linking spheres give a basis of
$H_{d-1}(U;\mathbb Z)$, with no other reduced homology in $U$.
For precision, excise the complement of the balls from the pair
$(\mathbb C^v,U)$. Relative homology becomes the direct sum of the
pairs $(B_a,B_a\setminus\{a\})$. Each summand is $\mathbb Z$ in
degree $d$ and zero otherwise. As $\mathbb C^v$ is contractible,
the boundary map in the relative long exact sequence identifies
these $N$ generators with the linking-sphere basis in $H_{d-1}(U)$.
Naturality and the projection retractions prove that the inclusion
map
\[
     \bigoplus_a H_{d-1}(W_a;\mathbb Z)
                 \longrightarrow H_{d-1}(A;\mathbb Z)
\]
is an isomorphism. This verifies the attaching map on the needed
homology group; a dimension count alone would not suffice.

Apply Mayer--Vietoris to the open cover $X(\mathbb C)=A\cup B$,
where $B=\coprod_a V_a$ and $A\cap B=\coprod_a W_a$.
Its degree-$d$ portion is
\[
 0\longrightarrow\bigoplus_a H_d(V_a;\mathbb Z)
    \longrightarrow H_d(X(\mathbb C);\mathbb Z)
    \longrightarrow\bigoplus_a H_{d-1}(W_a;\mathbb Z)
    \xrightarrow{\ \sim\ }H_{d-1}(A;\mathbb Z).
\]
Thus the local quadric sphere classes give a basis of
$H_d(X(\mathbb C);\mathbb Z)=\mathbb Z^N$.
All other reduced homology groups vanish. At the bottom of the
sequence, including when $d=2$, the map
$H_0(A\cap B)=\mathbb Z^N\to H_0(A)\oplus H_0(B)=\mathbb Z\oplus\mathbb Z^N$
is $e_a\mapsto(1,-e_a)$ and is injective. This accounts for the
potential degree-one boundary term and proves connectivity.

Simple connectivity also requires a check. Add the $V_a$ one at a
time to $A$. They are simply connected and intersect the preceding
union in the path-connected $W_a$, since the $B_a$ are disjoint.
Van Kampen consequently expresses the fundamental group as
$\pi_1(A)$ modulo the normal subgroup generated by the images of
all $\pi_1(W_a)$, with paths to a common basepoint understood.
Projection identifies this quotient with the quotient obtained by
adding the contractible $B_a$ to $U$ in the base. The latter union
is $\mathbb C^v$, whose fundamental group is trivial. Thus
$X(\mathbb C)$ is simply connected. This argument treats $v=1$
without replacing its nontrivial puncture loops by a homology-only
argument.

The spaces involved have CW homotopy type. One may see this directly
from this open cover: $U$ deformation retracts onto a large
rectangular box with the interiors of $N$ small rectangular boxes
removed, by pushing radially away from each puncture and inward
from infinity. Subdividing a finite rectangular grid triangulates
that region. The $V_a$ and $W_a$ have the sphere homotopy types just
proved. The double mapping cylinder for the two-set open cover is
homotopy equivalent to $X(\mathbb C)$; replacing its vertex spaces
by these CW models and its attaching maps by cellular
approximations yields a CW model. Homotopy invariance of this gluing
and the open-cover equivalence are Hatcher, Propositions4G.1--4G.2,
pp.458--459; cellular approximation is Theorem4.8, p.349.

Since the space is simply connected and its homology below degree
$d$ vanishes, successive applications of Hurewicz give
$\pi_i(X(\mathbb C))=0$ for $i<d$ and
$\pi_d(X(\mathbb C))\simeq H_d(X(\mathbb C);\mathbb Z)$.
Represent the displayed $N$ homology basis elements by based sphere
maps and form their wedge. This map from $\bigvee_N S^d$ induces an
isomorphism on all integral homology groups. Both spaces are simply
connected and have CW homotopy type, so the homology form of
Whitehead's theorem makes the map a homotopy equivalence.
The cohomology-ring assertion follows: a product of positive-degree
classes would lie in degree $4v$, whose group vanishes.

If $N=0$, then $U=\mathbb C^v$ and the first explicit deformation
retraction already makes $X(\mathbb C)$ contractible. There is in
fact the stronger arithmetic isomorphism claimed in the statement.
The integral Boolean quotient in NN-C027 is the product of $N$
copies of $\mathbb Z$. If $N=0$ it is the zero ring, so the ideal
$(f_1,\ldots,f_h)$ equals the whole of $R=\mathbb Z[x]$.
Choose a polynomial column $a=(a_1,\ldots,a_h)$ with $f\cdot a=1$.
For a point $(x,t,y)$ of $X$, put
\[
                         z=y+a(2t-t^2).
\]
Then $f\cdot z=t^2-t+2t-t^2=t$. Conversely, from arbitrary
$(x,z)\in\mathbb A^{v+h}_{\mathbb Z}$ set
\[
             t=f(x)\cdot z,\qquad y=z+a(x)(t^2-2t).
\]
These polynomial formulas are inverse over $\mathbb Z$, and
$f\cdot y=t^2-t$ proves that the second lands in $X$.
Thus $X\simeq\mathbb A^{v+h}_{\mathbb Z}$.
If $N>0$, its nonzero degree-$2v$ integral cohomology forbids
contractibility and forbids an isomorphism with complex affine
space. Base change then also forbids the integral isomorphism.
This proves all the equivalences.
\end{proof}

\paragraph{Exact topological source interfaces.}
Allen Hatcher, \emph{Algebraic Topology}, author text
\url{https://pi.math.cornell.edu/~hatcher/AT/AT.pdf}:
Theorem1.20, p.43 (van Kampen);
Theorem2.20, p.119 (excision);
Section2.2, ``Mayer--Vietoris Sequences'', pp.149--150 (unnumbered
exact sequence and its inclusion maps);
Theorem4.32, pp.366--367 (Hurewicz);
Corollary4.33, pp.367--368 (homology Whitehead);
Theorem4.8, p.349 and Propositions4G.1--4G.2, pp.458--459
(CW replacement and open-cover gluing).
The cited statements and the local proof passages used here were
read in the author PDF. The explicit affine-quadric retraction,
local analytic reduction, linking basis and arithmetic polynomial
isomorphism are proved above.

\paragraph{Effectivity and complexity scope.}
On the encoded family of NN-C027, $N$ is the satisfying-assignment
count. The rank of the single positive-degree integral cohomology
group, and the number of spheres in this homotopy type, are therefore
counting-complete quantities; their parity is parity-counting-complete.
Complex contractibility, isomorphism to complex affine space, and
isomorphism to integral affine space are each coNP-complete on this
encoded family, because each is equivalent to unsatisfiability.
Their complementary questions are NP-complete. Every scheme under
discussion is affine; the distinction is whether it is isomorphic
to affine space. The polynomial identity $f\cdot a=1$ and the
resulting isomorphism need not have a polynomial-size description,
so this construction supplies no short coNP certificate. E evaluates
these decisions through the original Boolean formula, while A and M
are not needed for the homotopy or isomorphism statements.
No motivic or $\mathbb A^1$-homotopy equivalence is asserted by this
complex-topological proof.
